\documentclass{article}
\usepackage{amsmath,amssymb}
\usepackage{xurl}
\usepackage[hidelinks]{hyperref}
\setlength{\emergencystretch}{2em}
% Shared commands for notes in docs/paper-gaps/.

\DeclareUnicodeCharacter{2080}{\ifmmode{}_{0}\else\textsubscript{0}\fi}
\DeclareUnicodeCharacter{2081}{\ifmmode{}_{1}\else\textsubscript{1}\fi}
\DeclareUnicodeCharacter{2082}{\ifmmode{}_{2}\else\textsubscript{2}\fi}
\DeclareUnicodeCharacter{2099}{\ifmmode{}_{n}\else\textsubscript{n}\fi}

\newcommand{\Lean}{\textsc{Lean}}
\ExplSyntaxOn
\NewDocumentCommand{\leanid}{m}{
  \ifmmode
    \textsf{\detokenize{#1}}
  \else
    \group_begin:
    \urlstyle{sf}
    \tl_set:Nn \l_tmpa_tl {#1}
    \tl_replace_all:Nnn \l_tmpa_tl {\_} {_}
    \exp_args:NV \path \l_tmpa_tl
    \group_end:
  \fi
}
\ExplSyntaxOff
\providecommand{\tr}{\operatorname{tr}}
\newcommand{\tnleanIssueBase}{https://github.com/LionSR/TNLean/issues/}
\newcommand{\tnleanPrBase}{https://github.com/LionSR/TNLean/pull/}
\newcommand{\tnleanDocBase}{https://sirui-lu.com/TNLean/docs/}
\newcommand{\ghissue}[1]{\href{\tnleanIssueBase#1}{GitHub issue \##1}}
\newcommand{\ghpr}[1]{\href{\tnleanPrBase#1}{PR \##1}}
\newcommand{\doclink}[1]{\href{\tnleanDocBase#1.html}{\path{#1}}}

% Dirac notation and the identity operator, as in blueprint/src/macros/common.tex.
% \providecommand keeps note-local definitions authoritative.
\usepackage{dsfont}
\providecommand{\ket}[1]{|#1\rangle}
\providecommand{\bra}[1]{\langle #1|}
\providecommand{\braket}[2]{\langle #1 | #2 \rangle}
\providecommand{\ketbra}[2]{|#1\rangle\!\langle #2|}
\providecommand{\Id}{\mathds{1}}
\providecommand{\one}{\mathds{1}}
\providecommand{\C}{\mathbb{C}}
\providecommand{\MN}[1]{M_{#1}(\C)}
\providecommand{\MD}{M_D(\C)}

% Verdict marker: \gapnote{<kind>}{<status>}. Kind is the policy
% classification (clarification, local-correction, scope-restriction,
% unfaithful, false-source, open-gap); status is open, wip (elimination
% underway), resolved, or historical. Machine-read by the site index;
% prints a verdict line.
\newcommand{\gapnote}[2]{\par\noindent\textbf{Verdict:} #1 (#2).\par}

\title{Weak-Hopf Parent Terms: Support Averaging and Ambient Completion}
\author{}
\date{2026-10-05}
\begin{document}
\maketitle
\gapnote{scope-restriction}{open}

\section{The source and the distinction}
Garre-Rubio, Lootens and Moln\'ar define the canonical local parent term
as $h=I-P$, where $P$ is the orthogonal projector onto the local
arbitrary-boundary matrix product state support. They average $h$ for a
group representation, then describe a weak-Hopf extension. Appendix B
explicitly specifies the weak-Hopf construction on $P$:
\[
    E(P)=\sum\rho(\Omega_{(1)})P\rho(T(\Omega_{(2)})),
    \qquad \rho=(\phi\otimes\phi)\circ\Delta.
\]
The normalization stated there is $E(I)=Q:=\rho(1)$, rather than
$E(I)=I$.\footnote{Exact source:
    \href{https://arxiv.org/abs/2203.12563v3}{arXiv:2203.12563v3},
    \path{Papers/2203.12563/REsubmission.tex}, lines 1276--1320 and 2463--2466.
    The canonical pulling-through objects are described in
    \href{https://arxiv.org/abs/2204.05940v1}{arXiv:2204.05940v1},
    Section 5.5, Theorem 5.4, and Section 6, equations (6.1)--(6.4).}
Consequently
\[
    E(I-P)=Q-E(P),\qquad I-E(P)=E(I-P)+(I-Q).
\]
The literal average of the parent and the ambient complement of the
averaged support are different constructions. If $E(P)$ is the
orthogonal projector onto the original support, as Appendix B states,
then $E(P)=P$ and the latter construction is the original parent term.
The nonzeroness of this complement additionally needs $P\ne I$.

\section{An elementary fusion-compatible test}
The following example is a mathematical derivation, independently checked
at the level of the displayed formulas; it is not a Lean weak-Hopf
algebra instance or a formal certificate of the complete weak-Hopf axioms.
Fix $r\ge2$, put $K=\mathbb C^r$, $V=K\otimes K$ and
$\mathcal A=\operatorname{End}(V)$. For matrix units
$e_{ij,kl}=|i,j\rangle\langle k,l|$, define
\[
    \Delta(e_{ij,kl})=\frac1r\sum_{a,b}e_{ia,kb}\otimes e_{aj,bl},\quad
    \epsilon(e_{ij,kl})=r\delta_{ij}\delta_{kl},\quad
    S(e_{ij,kl})=e_{lk,ji}.
\]
These formulas define a finite-dimensional $C^*$-weak Hopf algebra.
Indeed, with
\[
    W_2|i,j\rangle=\frac1{\sqrt r}\sum_a|i,a\rangle|a,j\rangle,
    \qquad |\Phi\rangle=\frac1{\sqrt r}\sum_a|a,a\rangle,
\]
one has $W_2^*W_2=I$ and $\Delta(x)=W_2xW_2^*$. Coassociativity
inserts the same two Bell pairs in either order. The weak-unit identity
is the equality of the product of the two disjoint Bell projections
with $\Delta^2(1)$. The counit and weak-counit identities follow by
contracting their Kronecker deltas. The counital maps are
\[
    \epsilon_t(e_{ij,kl})=\delta_{kl}(E_{ij}\otimes I),\qquad
    \epsilon_s(e_{ij,kl})=\delta_{ij}(I\otimes E_{kl}),
\]
and direct multiplication gives all three antipode identities.

The positive normalized integral and its pulling-through map are
\[
    \Omega=\frac1r\sum_{i,k}e_{ii,kk}=|\Phi\rangle\langle\Phi|,
    \qquad T=S.
\]
In particular, $\Delta(\Omega)=r^{-2}\sum_{i,k,a,b}
    e_{ia,kb}\otimes e_{ai,bk}$ is cocommutative and nondegenerate, and
$(1\otimes x)\Delta(\Omega)=(S(x)\otimes1)\Delta(\Omega)$.
The dual is again this weak Hopf algebra: for the dual basis,
$f_{ij,kl}\mapsto r^{-1}e_{ik,jl}$ is a star-preserving weak-Hopf
isomorphism. Its target counital subalgebra is the irreducible defining
module, and $W_2$ identifies its weak tensor square with itself. Thus
both representation categories are equivalent to $\mathrm{Vec}$, with
one simple tensor unit, rather than merely to a multifusion category.
The unique dual character equals $\Omega$, so the canonical character
construction gives exactly the displayed normalization.

Consider the injective matrix product state and operator tensors
\[
    A^{ij}=E_{ij}/\sqrt r,\qquad
    T^{ij,kl}=(E_{ij}\otimes E_{kl})/r.
\]
Their matrices span their respective full bond matrix algebras. Define
$W_n:V\to V^{\otimes n}$ by inserting $n-1$ normalized Bell pairs.
The arbitrary-boundary state support is $\operatorname{ran}W_n$.
For an operator boundary $B$, set
\[
    x(B)_{(i,j),(k,l)}=r^{-1}B_{(j,l),(i,k)}.
\]
The corresponding operator is $O^n_{T,B}=W_nx(B)W_n^*$ at every positive
length. This length-independent bijection of boundaries proves
arbitrary-boundary closedness. Its action on state boundary vectors
proves the length-independent compatibility condition as well.
For $B=I$, $x(B)=\Omega$ and the periodic operator is the rank-one
projection onto the periodic state. The fusion and action reductions
insert one normalized Bell pair; their associativity coefficients can
be chosen $F=L=1$.

At two sites the support is
\[
    P=Q=\Delta(1)=I_K\otimes|\Phi\rangle\langle\Phi|\otimes I_K.
\]
Writing $Y=W_2^*XW_2$ yields
\[
    E(X)=W_2\left(\frac1{r^2}\sum_{i,a,k,b}
    e_{ia,kb}Y e_{kb,ia}\right)W_2^*
    =\frac{\operatorname{Tr}(QX)}{r^2}Q.
\]
Thus $E(P)=P$, while $E(I-P)=0$ and
$\operatorname{rank}(I-P)=r^4-r^2>0$. This disproves unconditional
nonzeroness for the \emph{uncompleted average}. It is fully consistent
with Appendix B read as choosing the parent $I-E(P)$. No counterexample
to the phase-classification theorem is asserted.

\section[The local theorem and the remaining formalization scope]
 {The local theorem and the remaining\\formalization scope}
Let $\mathcal B$ be a star-closed family of matrices preserving a
subspace $G$. For $v\in G^\perp$ and $w\in G$,
\[
    \langle w,Bv\rangle=\langle B^*w,v\rangle=0.
\]
Hence both $G$ and $G^\perp$ are invariant, and their orthogonal
projections commute with every member of $\mathcal B$. Adjoining the
ambient identity does not change invariance.

The local parent theorem derives invariance from the source's
arbitrary-boundary compatibility, requires adjoint closure of the
local boundary range explicitly, and applies this orthogonal-reduction
argument to the canonical complementary projection.\footnote{
    \leanid{MPOTensor.IsBoundaryCompatible.groundSpaceES_mem_invtSubmodule},
    \leanid{MPOTensor.IsBoundaryCompatible.groundSpaceES_invariant_boundaryStarAlgebra},
    \leanid{MPOTensor.IsBoundaryCompatible.parentInteractionES_commute_mpoWithBoundary},
    and
    \leanid{MPOTensor.IsBoundaryCompatible.parentInteraction_commute_mpoWithBoundary},
    in \doclink{TNLean/MPS/Symmetry/BoundaryParentCommutation}.}
It does not assume the desired commutation, nor does it identify a
nonunital weak-Hopf average with an ambient complementary projection.

\section[Unit support and normalized boundary averages]
 {Unit support and normalized\\boundary averages}
There is a sufficient positive-length criterion for the missing identity
$Q_nP_n=P_n$. Let $T$ be an operator tensor and $A$ a state tensor,
compatible with one output state boundary chosen for all positive lengths.
Assume that $A$ is injective after blocking $L_0>0$ sites and that a fixed
operator boundary $U$ satisfies $O^{L_0}_{T,U}=I$. Write
$\Gamma_k(X)=\psi^k_{A,X}$ and $Q_n=O^n_{T,U}$. Compatibility gives
one state boundary $Z$ with
\[
    O^k_{T,U}\Gamma_k(X)=\Gamma_k(Z)
\]
for all $k>0$. At $k=L_0$ this becomes
$\Gamma_{L_0}(Z)=\Gamma_{L_0}(X)$, so injectivity forces $Z=X$.
Consequently $Q_n$ fixes every vector of $\mathcal S_A^n$ and
$Q_nP_n=P_n$ for every $n>0$. The target length can be smaller than
$L_0$, and no adjoint closure is needed for this one-sided identity.
One-site injectivity and $O^1_{T,U}=I$ are a special case.\footnote{
    The boundary-vector, whole-support, and canonical-projection versions,
    including one-site specializations, are proved in
    \doclink{TNLean/MPS/Symmetry/BoundaryUnitSupport}.}

Now also assume adjoint closure of the entire boundary range at the
chosen target length $n>0$. Orthogonal reduction gives
$P_nO^n_{T,Y}=O^n_{T,Y}P_n$ for every boundary $Y$. Let
$(X_j,Y_j)_{j\in J}$ be finite actual boundary families with
\[
    \sum_{j\in J}O^n_{T,X_j}O^n_{T,Y_j}=Q_n,
    \qquad E_n(B)=\sum_{j\in J}O^n_{T,X_j}BO^n_{T,Y_j}.
\]
This is the product normalization used in Appendix B. The derived
commutation and unit-support identities imply
\[
    E_n(P_n)=\sum_{j\in J}(O^n_{T,X_j}O^n_{T,Y_j})P_n=Q_nP_n=P_n.
\]
It follows that
\[
    E_n(I-P_n)=Q_n-P_n,\qquad I-E_n(P_n)=I-P_n.
\]
The ambient complement is nonzero whenever $d^n>D^2$, by the dimension
bound $\dim\mathcal S_A^n\leq D^2$. The literal parent average can
still vanish when $Q_n=P_n$. The finite index set may be empty and the dimensions may be zero
whenever the displayed assumptions hold. No reverse product
normalization, paired-adjoint identity, or assumed value of $E_n(P_n)$
is required.\footnote{
    The commutation, support-average, literal parent-average, ambient
    complement, and nonzero-complement results are proved in
    \doclink{TNLean/MPS/Symmetry/BoundaryProjectionAverage}.}

\section{Remaining weak-Hopf construction}
The source-labelled weak-Hopf assertion remains separate. Its proof requires
the $C^*$-weak-Hopf representation, its boundary adjoint closure, a represented
unit satisfying the unit-support condition, and normalized finite boundary
families from the canonical integral. The positive injectivity and identity
hypotheses above are not derived from general boundary compatibility.

The periodic-parent results supply tensor cuts and cyclic-window transport
under explicit local adjoint closure. Wrapping windows require the global
boundary to commute with each tensor letter. Summand selectors of assembled
block tensors satisfy this condition and extract the labelled periodic
operators. These results supply neither the missing weak-Hopf realization
nor its integral normalization.
\end{document}
